\documentclass[aspectratio=169]{beamer}
\usetheme{Madrid}
\usepackage[utf8]{inputenc}
\usepackage[T1]{fontenc}
\usepackage{tikz}
\usepackage{amssymb}
\usepackage{graphicx}
\setbeamertemplate{navigation symbols}{}
\setbeamercolor{title}{fg=white,bg=orange!85!red}
\setbeamercolor{frametitle}{fg=white,bg=orange!85!red}
\setbeamercolor{block title}{bg=orange!85!red!20,fg=black}
\setbeamercolor{block body}{bg=orange!85!red!8,fg=black}
\setbeamercolor{item}{fg=orange!85!red}
\setbeamercolor{itemize item}{fg=orange!85!red}
\setbeamertemplate{itemize item}[triangle]

\title{What Is a Website?}
\subtitle{Client/server, URLs, real browsers — no code yet, just the real parts}
\author{ByteForge --- Grade 8 --- App Builder}
\date{}

\begin{document}

\begin{frame}
    \titlepage
\end{frame}

\begin{frame}{Today's Plan}
    \begin{itemize}
        \item Client \& Server
        \item Real Browsers \& Devices
        \item[$\bigstar$] BOSS --- Inspect a Real Page
    \end{itemize}
\end{frame}

\begin{frame}{Browser vs. server — two genuinely different jobs}
    \large
    \begin{itemize}
        \item A browser is the customer that asks for a page; a server is the kitchen that actually stores the files and sends them back
        \item A real, physical machine somewhere, not just an app on your device.
    \end{itemize}
\end{frame}


\begin{frame}{A URL's parts each do a different job}
    \large
    \begin{itemize}
        \item \texttt{https://} is the protocol
        \item HOW to connect; \texttt{wikipedia.org} is the domain
        \item The site's real address. Neither one is optional or interchangeable.
    \end{itemize}
\end{frame}


\begin{frame}{JavaScript and Java are genuinely different languages}
    \large
    \begin{itemize}
        \item Despite the similar name, JavaScript (a web scripting language) and Java (a completely separate language) share nothing but a marketing-era name collision from the 1990s
        \item A real, common point of confusion worth clearing up early.
    \end{itemize}
\end{frame}


\begin{frame}{HTML/CSS/JS as skeleton, paint, electricity}
    \large
    \begin{itemize}
        \item HTML is a building's skeleton (structure), CSS is the paint (style), JS is the electricity (things that move and respond)
        \item Three separate real jobs, never one file doing all three.
    \end{itemize}
\end{frame}


\begin{frame}{BOSS — Page Source vs. Inspect, a real distinction}
    \large
    \begin{itemize}
        \item Page Source shows the ORIGINAL html the server sent, unchanged; Inspect shows the page's LIVE current state, which JavaScript can have already changed
        \item Genuinely different views of the same page.
    \end{itemize}
\end{frame}

\begin{frame}{Real-World Examples}
    \large
    \begin{itemize}
        \item Every time a phone or laptop loads a webpage, it is genuinely acting as a client asking a real server (often thousands of miles away) for files
        \item The same client/server split every professional web developer works with daily, not a simplification invented for class.
    \end{itemize}
\end{frame}


\begin{frame}{Let Them Discover}
    \Large
    Right-click → Inspect on three real websites you use. Find one heading, one button, one link, and one image in each site's Inspect view — notice the same handful of tags underneath every single site, no matter how different they look.
\end{frame}

\begin{frame}{Your Turn}
    \Large
    Open the student map, find \textbf{What Is a Website?}, and work through it
    hands-on --- everything you just saw, they now do themselves.
\end{frame}

\end{document}
